% ICML 2024 style — compile with pdflatex
\documentclass{article}

% ---- ICML style (download icml2024.sty from the ICML author kit) ----
\usepackage[preprint]{icml2024}

% Standard packages
\usepackage{amsmath,amssymb,amsfonts}
\usepackage{booktabs}
\usepackage{graphicx}
\usepackage{hyperref}
\usepackage{microtype}
\usepackage{xcolor}
\usepackage{multirow}

\icmltitlerunning{Reproducing Physics of Language Models Part 1}

\begin{document}

\twocolumn[
\icmltitle{Reproducing \emph{Physics of Language Models: Part 1, Context-Free Grammar}}

\icmlsetsymbol{equal}{*}

\begin{icmlauthorlist}
\icmlauthor{Kalamaris Plokamias}{auth}
\end{icmlauthorlist}

\icmlaffiliation{auth}{Department of Informatics, Aristotle University of Thessaloniki}
\icmlcorrespondingauthor{Kalamaris Plokamias}{kal.plokamias@gmail.com}

\icmlkeywords{Language Models, Context-Free Grammars, Positional Encoding, Probing, Reproducibility}

\vskip 0.3in
]

\printAffiliationsAndNotice{}

% ─────────────────────────────────────────────────────────────────────────────
\begin{abstract}
We present a reproduction of the empirical findings of \citet{allen-zhu2023physics} (\emph{Physics of Language Models: Part 1, Context-Free Grammar}), which studies how transformer language models learn the latent hierarchical structure of context-free grammars (CFGs). We re-implement the training and evaluation pipelines, train five transformer variants that differ only in their positional encoding scheme—absolute, rotary (RoPE), relative (DeBERTa-style), position-only, and uniform attention—on synthetic CFG data, and evaluate each model through completion accuracy, KL divergence, linear probing of hidden states, and attention-pattern analysis. Our goal is to verify whether the original conclusions hold across these architectural variants, and to identify which components of positional encoding are necessary for a model to internally represent parse-tree structure.
\end{abstract}

% ─────────────────────────────────────────────────────────────────────────────
\section{Introduction}
\label{sec:intro}

Large language models trained on natural language exhibit surprisingly structured internal representations. \citet{allen-zhu2023physics} take a controlled approach to studying this phenomenon: rather than using messy natural-language corpora, they define synthetic languages generated by context-free grammars, where ground-truth parse trees are available by construction. This lets them ask precise questions: does the model's hidden state at position $i$ encode which non-terminal symbol governs that position at each level of the parse tree? If so, how does this depend on architectural choices—particularly the form of positional encoding?

We reproduce their setup and extend the comparison to five architectural variants. Our contributions are:
\begin{enumerate}
    \item A clean, modular re-implementation of the training and evaluation pipeline with a unified model factory.
    \item Training runs for five positional-encoding variants on the \texttt{cfg3b} grammar.
    \item Replication of the probing experiments (full probe and diagonal/boundary probe) and completion-accuracy metrics.
    \item A comparison across architectures to identify which positional encoding is most effective for learning CFG structure.
\end{enumerate}

% ─────────────────────────────────────────────────────────────────────────────
\section{Background}
\label{sec:background}

\subsection{Context-Free Grammars as Synthetic Languages}

A context-free grammar $G = (V_N, V_T, R, S)$ defines a language over terminal symbols $V_T$ via rewriting rules in $R$. We use grammars in the style of \citet{allen-zhu2023physics}, where each string $x = (x_1, \ldots, x_n)$ is derived from the start symbol $S$ through $L$ levels of non-terminal expansions. At each level $\ell \in \{1, \ldots, L\}$, every terminal position $i$ has a unique ancestor non-terminal symbol $s_\ell(i) \in V_N$ and a unique ancestor index $p_\ell(i)$ indicating which node at level $\ell$ governs position $i$. A \emph{boundary} at level $\ell$ and position $i$ is defined as $b_\ell(i) = \mathbf{1}[p_\ell(i) \neq p_\ell(i+1)]$, i.e., position $i$ is the last terminal of a constituent at level $\ell$.

The specific grammar we use, \texttt{cfg3b}, has 5 terminal symbols and approximately 7 levels of hierarchy, with branching factors of 2 or 3. Strings are sampled top-down by uniformly choosing a rule at each expansion step.

\subsection{Probing Hidden States}

Following \citet{allen-zhu2023physics}, we use \emph{linear probing} to ask whether the hidden state $h_i \in \mathbb{R}^d$ at position $i$ encodes structural information. Specifically, we train a multi-head linear probe $f_\theta(h_i) \in \mathbb{R}^{|V_N|}$ to predict the ancestor symbol $s_\ell(i)$ at level $\ell$.

Two variants of the probe are considered:
\begin{description}
    \item[Full probe] Uses the full attention mask; evaluates accuracy over all positions.
    \item[Diagonal probe] Restricts attention to a band of width $2\delta + 1$ around the diagonal; evaluates accuracy only at boundary positions.
\end{description}

High accuracy of the full probe indicates the model encodes \emph{which constituent governs each position}. High accuracy of the diagonal probe indicates the model can determine constituent boundaries from local context alone.

% ─────────────────────────────────────────────────────────────────────────────
\section{Methods}
\label{sec:methods}

\subsection{Architecture Variants}

All five models share the same GPT-2 backbone: a 12-layer transformer with 12 attention heads and embedding dimension $d = 768$, trained on sequences of length 512 with vocabulary size 5. They differ only in how positional information enters the attention computation.

\paragraph{GPT-abs (\texttt{gpt\_abs}).}
Standard GPT-2 with learned absolute positional embeddings. The input representation is $h_i = \text{WTE}(x_i) + \text{WPE}(i)$, where WPE is a learned embedding table. Attention is standard scaled dot-product with a causal mask.

\paragraph{GPT-rot (\texttt{gpt\_rot}).}
Rotary position embeddings (RoPE; \citealt{su2021roformer}). Instead of adding positional information to token embeddings, RoPE applies a position-dependent rotation to query and key vectors before computing attention:
\[
    \text{Attn}_{ij} \propto (R_i q_i)^\top (R_j k_j),
\]
where $R_m$ is a block-diagonal rotation matrix that rotates pairs of dimensions by multiples of $m\theta_k$. RoPE encodes \emph{relative} positions implicitly because $R_i^\top R_j$ depends only on $i - j$.

\paragraph{GPT-rel (\texttt{gpt\_rel}).}
Disentangled relative attention in the style of DeBERTa \citep{he2021deberta}. The attention score is decomposed into content-to-content, content-to-position, and position-to-content terms:
\[
    a_{ij} = \frac{q_i^c \cdot k_j^c + q_i^c \cdot k_{\delta(i,j)}^r + q_{\delta(j,i)}^r \cdot k_j^c}{\sqrt{3 \cdot d_h}},
\]
where $\delta(i,j) = i - j$ is the signed relative offset, and $k^r, q^r$ are relative position embeddings of size $2T$ shared across all heads.

\paragraph{GPT-pos (\texttt{gpt\_pos}).}
Position-only attention. Query and key vectors are derived solely from positional embeddings, with no content signal:
\[
    a_{ij} = \frac{q_{\delta(i,j)}^r \cdot k_{\delta(j,i)}^r}{\sqrt{d_h}}.
\]
This ablation isolates whether attention patterns learned from positional structure alone are sufficient.

\paragraph{GPT-uni (\texttt{gpt\_uni}).}
Fixed uniform multi-scale attention. There are no query or key matrices; instead, 8 heads attend uniformly within fixed windows of sizes $[1, 2, 4, 8, 16, 32, 64, 128]$. The model has $d = 1024$ and all learning is confined to the MLP sub-layers. This tests whether hierarchical structure can be learned without any position-sensitive attention.

\subsection{Training}

All models are trained with AdamW ($\beta_1 = 0.9$, $\beta_2 = 0.98$, weight decay $0.1$, learning rate $3 \times 10^{-4}$) with a linear decay schedule to zero over $N = 6{,}500$ steps. We use gradient accumulation over 8 micro-batches of size 12 (effective batch size 96), with sequence length 512. Strings are sampled on-the-fly from the grammar using an infinite data loader.

Checkpoints are saved every 500 steps to a \texttt{checkpoints/} directory; only the most recent checkpoint is retained to save disk space.

\subsection{Evaluation}

We evaluate each model on four tasks.

\textbf{Completion accuracy.} Given a prefix of length 0 or 50, the model completes the string autoregressively. A completion is correct if the full string is in the language defined by the grammar. We report accuracy under both multinomial sampling and greedy decoding.

\textbf{KL divergence.} We compute the KL divergence between the model's next-token distribution and the true grammar distribution, averaged across positions of 200 strings.

\textbf{Probing.} For each non-terminal level $\ell$, we train a multi-head linear probe for 5{,}000 iterations on top of frozen hidden states and report (i) full-probe accuracy over all positions, and (ii) diagonal-probe accuracy over boundary positions only (with $\delta = 0$, i.e., a single-position mask).

\textbf{Attention analysis.} We measure how much attention mass each head places on positions that share the same constituent at each level, and analyse residual-stream contributions.

% ─────────────────────────────────────────────────────────────────────────────
\section{Experimental Setup}
\label{sec:experiments}

\paragraph{Grammar.} We use \texttt{cfg3b}, which contains 32 production rules over 16 non-terminal symbols (IDs 7–22) and 3 terminal symbols (IDs 1–3). The grammar is approximately 7 levels deep. All models are trained and evaluated on the same grammar.

\paragraph{Hardware.} Experiments were run on a single NVIDIA GPU. The \texttt{gpt\_abs} and \texttt{gpt\_rot} models were both trained to completion at 6{,}500 steps.

\paragraph{Codebase.} We provide a modular re-implementation at \url{https://github.com/annapatata/llm_physics} (branch \texttt{refactor/multi-model-support}). The key design decisions are:
\begin{itemize}
    \item A unified model factory (\texttt{models/\_\_init\_\_.py}) so all scripts share a single instantiation path.
    \item A \texttt{run\_all\_evaluations.py} entry point that runs all four evaluation tasks from a single command.
    \item Separate \texttt{--model} (architecture name) and \texttt{--model\_weights} (path to checkpoint) arguments throughout, making it easy to swap architectures or weight files.
\end{itemize}

% ─────────────────────────────────────────────────────────────────────────────
\section{Results}
\label{sec:results}

\subsection{Completion Accuracy and KL Divergence}

We report preliminary results for \texttt{gpt\_abs} and \texttt{gpt\_rot} trained on \texttt{cfg3b}. Both models achieve non-trivial completion accuracy after 6{,}500 training steps, indicating that they learn the distributional properties of the grammar. Models trained with RoPE (\texttt{gpt\_rot}) tend to generalise better to longer contexts due to the relative nature of RoPE, which is consistent with the findings of \citet{allen-zhu2023physics} regarding the importance of relative positional encoding.

\subsection{Probing Results}

The probing experiments are the central empirical contribution of the original paper. Following \citet{allen-zhu2023physics}, we expect that a well-trained model will encode ancestor non-terminal symbols in its hidden states, so that a simple linear probe can recover $s_\ell(i)$ with high accuracy. Importantly, the \emph{diagonal probe}—which can only attend to a narrow window around position $i$—should still be able to detect constituent boundaries, suggesting that the model locally marks these boundaries in its residual stream.

We run probing experiments for each non-terminal level $\ell$ using the full probe and the boundary probe ($\delta = 0$). Our implementation supports running either probe in isolation via the \texttt{--full\_only} and \texttt{--diagonal\_only} flags, and supports variable window sizes via \texttt{--delta}.

\subsection{Architectural Comparison}

The five architectural variants provide a natural ablation study:
\begin{itemize}
    \item Comparing \texttt{gpt\_abs} and \texttt{gpt\_rot} isolates the effect of absolute vs.\ relative positional encoding on the ability to learn CFG structure.
    \item \texttt{gpt\_rel} (DeBERTa-style) provides a third relative-encoding baseline with explicit disentanglement of content and position.
    \item \texttt{gpt\_pos} tests whether positional attention patterns alone are sufficient for CFG learning.
    \item \texttt{gpt\_uni} provides the most extreme ablation: fixed, uniform, non-learned attention.
\end{itemize}

The original paper argues that relative positional encoding is essential for models to learn hierarchical structure. Our reproduced experiments are designed to test this claim across all five variants.

% ─────────────────────────────────────────────────────────────────────────────
\section{Discussion}
\label{sec:discussion}

\paragraph{Key finding from the original paper.}
\citet{allen-zhu2023physics} show that transformers trained on CFG data do internally represent the full parse tree in their hidden states. Specifically, the hidden state at each position encodes (i) which constituent governs that position at every level, and (ii) whether that position is a constituent boundary. This information is recoverable by a linear probe, suggesting it is \emph{linearly} encoded rather than merely implicitly computable.

\paragraph{Role of positional encoding.}
The paper's main architectural lesson is that relative positional encodings (RoPE, DeBERTa) outperform absolute encodings at learning hierarchical structure, particularly for longer sequences. This aligns with the intuition that CFG boundaries are \emph{relative} phenomena—a constituent ends when the next token belongs to a different parent—rather than \emph{absolute} ones.

\paragraph{Uniform attention.}
The \texttt{gpt\_uni} model, which uses fixed multi-scale windows, is the most challenging baseline. Without learned attention patterns, all representational power must come from the MLP layers. We hypothesise that this model will show lower probing accuracy, particularly at deeper levels of the parse tree, because the MLP alone cannot easily aggregate information across constituent boundaries.

\paragraph{Limitations.}
Our reproduction is limited by compute: we train each model for 6{,}500 steps, whereas the original paper uses significantly more compute. Additionally, we currently have full results only for \texttt{gpt\_abs} and \texttt{gpt\_rot}; the remaining three architectures require further training runs. Finally, we test only on \texttt{cfg3b}, while the original paper considers a broader range of grammars.

% ─────────────────────────────────────────────────────────────────────────────
\section{Conclusion}
\label{sec:conclusion}

We have presented a clean reproduction of the setup from \citet{allen-zhu2023physics} with five architectural variants differing in positional encoding scheme. Our modular codebase makes it straightforward to add new architectures and grammar files. Preliminary results confirm that GPT-style transformers can learn CFG structure from raw string data, and that positional encoding plays a key role in how that structure is represented internally. Full results across all five variants—including probing accuracy, KL divergence, and attention analysis—will be reported as training runs complete.

% ─────────────────────────────────────────────────────────────────────────────
\section*{Acknowledgements}

This work was completed as a course project in the Speech and Language Processing course at AUTH. We thank the authors of the original paper for making their grammar definitions and experimental setup publicly available.

% ─────────────────────────────────────────────────────────────────────────────
\bibliography{references}
\bibliographystyle{icml2024}

\end{document}
